% Part: model-theory
% Chapter: models-of-arithmetic
% Section: standard-models

\documentclass[../../../include/open-logic-section]{subfiles}

\begin{document}

\olfileid[tr]{mod}{mar}{stm}
\olsection{Aritmetiğin Standart Modelleri}

Aritmetik dili~$\Lang{L_A}$ açıkça sayılar, özellikle doğal sayılar
hakkında olmak üzere tasarlanmıştır. Bu nedenle ``standart''
model~$\Struct N$ özeldir: Üzerinde konuşmak istediğimiz model odur.
Fakat mantıkta çoğu kez yalnızca yapısal özelliklerle ilgileniriz ve
izomorf olan herhangi iki !!{structure}s bu özellikleri paylaşır.
Dolayısıyla biraz daha geniş davranıp $\Struct N$ ile izomorf her
!!{structure} ifadesini ``standart'' sayabiliriz.

\begin{defn}
$\Lang{L_A}$ için bir !!{structure}, $\Struct N$ ile izomorfsa
\emph{standarttır}.
\end{defn}

\begin{prop}
\ollabel{prop:standard-domain} !!a{structure}~$\Struct M$ standartsa
evreni standart sayı terimlerinin değerleri kümesidir; yani
\[
\Domain M=\Setabs{\Value{\num n}{M}}{n\in\Nat}.
\]
\end{prop}

\begin{proof}
Her $\Value{\num n}{M}$ değerinin $\Domain M$ içinde olduğu açıktır.
Yalnızca her $x\in\Domain M$ için $x=\Value{\num n}{M}$ olacak bir
$n$ bulunduğunu göstermeliyiz. $\Struct M$ standart olduğundan
$\Struct N$ ile izomorftur. $g\colon\Nat\to\Domain M$ bir izomorfizma
olsun. O zaman
$g(n)=g(\Value{\num n}{N})=\Value{\num n}{M}$. $g$ !!{surjective}
olduğundan her $x\in\Domain M$ için $g(n)=x$ olacak bir $n\in\Nat$
vardır.
\end{proof}

\begin{explain}
$\Lang{L_A}$ için bir $\Struct M$ yapısı standartsa !!{domain}
içindeki bütün elemanlar standart $\num0,\num1,\num2,\dots$ sayı
terimleriyle, yani $\Obj0,\Obj0',\Obj0'',\dots$ terimleriyle
adlandırılabilir. Elbette bu, $\Domain M$ kümesinin !!{element}s
gerçekten sayılardır demek değildir; yalnızca onları $\Domain N$
içindeki sayıları seçtiğimiz biçimde seçebildiğimiz anlamına gelir.
\end{explain}

\begin{prob}
\olref[mod][mar][stm]{prop:standard-domain} önermesinin tersinin yanlış
olduğunu gösterin; yani
$\Domain M=\Setabs{\Value{\num n}{M}}{n\in\Nat}$ olduğu halde
$\Struct N$ ile izomorf olmayan !!a{structure}~$\Struct M$ örneği verin.
\end{prob}

\begin{prop}
\ollabel{prop:thq-standard}
$\Sat{M}{\Th Q}$ ve
$\Domain M=\Setabs{\Value{\num n}{M}}{n\in\Nat}$ ise $\Struct M$
standarttır.
\end{prop}

\begin{proof}
$\Struct M$ yapısının $\Struct N$ ile izomorf olduğunu göstermeliyiz.
$g(n)=\Value{\num n}{M}$ ile tanımlanan
$g\colon\Nat\to\Domain M$ fonksiyonunu ele alalım. Varsayım gereği
$g$ !!{surjective}. Ayrıca !!{injective}: $n\neq m$ olduğunda
$\Th Q\Proves\eq/[\num n][\num m]$. $\Sat{M}{\Th Q}$ olduğundan
$n\neq m$ ise $\Sat{M}{\eq/[\num n][\num m]}$. Dolayısıyla
$\Value{\num n}{M}\neq\Value{\num m}{M}$, yani $g(n)\neq g(m)$.

$g$'nin bir izomorfizma olduğunu da doğrulamalıyız.
\begin{enumerate}
\item $\Assign{\Obj0}{N}=0$ olduğundan
  $g(\Assign{\Obj0}{N})=g(0)$. $g$ tanımı gereği
  $g(0)=\Value{\num0}{M}$. Fakat $\num0$, $\Obj0$ teriminin
  kendisidir ve !!a{constant} olan bir terimin değeri, !!{structure}
  tarafından bu !!{constant} ifadesine atanan nesnedir; yani
  $\Value{\Obj0}{M}=\Assign{\Obj0}{M}$. Böylece gerektiği gibi
  $g(\Assign{\Obj0}{N})=\Assign{\Obj0}{M}$.
\item $\Struct N$ içindeki $\prime$, $\Nat$ üzerindeki ardıl fonksiyonu
  olduğundan $g(\Assign{\prime}{N}(n))=g(n+1)$. $g$ tanımı gereği
  $g(n+1)=\Value{\num{n+1}}{M}$. Ancak $\num{n+1}$ ile $\num n'$
  aynı terimdir; dolayısıyla
  $\Value{\num{n+1}}{M}=\Value{\num n'}{M}$. Değer fonksiyonunun
  tanımı gereği bu,
  $\Assign{\prime}{M}(\Value{\num n}{M})$ değeridir.
  $\Value{\num n}{M}=g(n)$ olduğundan
  $g(\Assign{\prime}{N}(n))=\Assign{\prime}{M}(g(n))$.
\item $\Struct N$ içindeki $+$, $\Nat$ üzerindeki toplama fonksiyonu
  olduğundan $g(\Assign{+}{N}(n,m))=g(n+m)$. $g$ tanımı gereği
  $g(n+m)=\Value{\num{n+m}}{M}$. Fakat
  $\Th Q\Proves\num{n+m}=(\num n+\num m)$; dolayısıyla
  $\Value{\num{n+m}}{M}=\Value{\num n+\num m}{M}$. Değer
  fonksiyonunun tanımı gereği bu,
  $\Assign{+}{M}(\Value{\num n}{M},\Value{\num m}{M})$ değeridir.
  $\Value{\num n}{M}=g(n)$ ve $\Value{\num m}{M}=g(m)$ olduğundan
  $g(\Assign{+}{N}(n,m))=\Assign{+}{M}(g(n),g(m))$.
\item $g(\Assign{\times}{N}(n,m))=
  \Assign{\times}{M}(g(n),g(m))$: Alıştırma.
\item $\tuple{n,m}\in\Assign{<}{N}$ ancak ve ancak $n<m$. $n<m$ ise
  $\Th Q\Proves\num n<\num m$ ve $\Sat{M}{\num n<\num m}$.
  Dolayısıyla
  $\tuple{\Value{\num n}{M},\Value{\num m}{M}}\in\Assign{<}{M}$,
  yani $\tuple{g(n),g(m)}\in\Assign{<}{M}$. $n\not<m$ ise
  $\Th Q\Proves\lnot\num n<\num m$ ve dolayısıyla
  $\Sat/{M}{\num n<\num m}$. Böylece
  $\tuple{g(n),g(m)}\notin\Assign{<}{M}$. Birlikte,
  $\tuple{n,m}\in\Assign{<}{N}$ ancak ve ancak
  $\tuple{g(n),g(m)}\in\Assign{<}{M}$ elde edilir.
\end{enumerate}
\end{proof}

\begin{explain}
$g$, $\Nat$ kümesinden $\Lang{L_A}$ için başka herhangi bir
!!{structure}~$\Struct M$ evrenine bir eşleme tanımlamanın en açık
yoludur; çünkü her böyle $\Struct M$, $\num0,\num1,\num2,\dots$
tarafından adlandırılan !!{element}s içerir. Dolayısıyla $\Struct M$,
$\num n$ terimleri hakkındaki en azından bazı temel önermeleri
$\Struct N$ ile aynı biçimde doğru kılıyorsa ve $g$ de bire bir örtense
$g$'nin bir izomorfizmaya dönüşmesi şaşırtıcı değildir. Hatta
$\Domain M$, $\num n$ terimlerinin adlandırdıkları dışında hiçbir
!!{element}s içermiyorsa tek izomorfizma odur.
\end{explain}

\begin{prop}
\ollabel{prop:thq-unique-iso} $\Struct M$ standartsa
\olref{prop:thq-standard} ispatındaki $g$, $\Struct N$ yapısından
$\Struct M$ yapısına tek izomorfizmadır.
\end{prop}

\begin{proof}
$h\colon\Nat\to\Domain M$, $\Struct N$ ile $\Struct M$ arasında bir
izomorfizma olsun. $n$ üzerinde tümevarımla $g=h$ olduğunu gösterelim.
$n=0$ ise $g$ tanımı gereği $g(0)=\Assign{\Obj0}{M}$. $h$ bir
izomorfizma olduğundan
$h(0)=h(\Assign{\Obj0}{N})=\Assign{\Obj0}{M}$; dolayısıyla
$g(0)=h(0)$.

Şimdi $n+1$ durumunu ele alalım:
\begin{align*}
  g(n+1) & = \Value{\num{n+1}}{M} && \text{$g$ tanımı gereği}\\
  & = \Value{\num n'}{M} && \text{$\num{n+1}\ident\num n'$ olduğundan}\\
  & = \Assign{\prime}{M}(\Value{\num n}{M})
    && \text{$\Value{t'}{M}$ tanımı gereği}\\
  & = \Assign{\prime}{M}(g(n)) && \text{$g$ tanımı gereği}\\
  & = \Assign{\prime}{M}(h(n)) && \text{tümevarım varsayımı gereği}\\
  & = h(\Assign{\prime}{N}(n)) && \text{$h$ izomorfizma olduğundan}\\
  & = h(n+1).
\end{align*}
\end{proof}

\begin{explain}
Her !!{denumerable} $M$ kümesi için $\Nat$ ile $M$ arasında
!!a{bijection} bulunduğundan, böyle her $M$ kümesi bir standart
$\Struct M$ modelinin !!{domain} olmaya adaydır. Aslında
$z\in M$ nesnesini ve uygun bir $s$ fonksiyonunu sırasıyla
$\Assign{\Obj0}{M}$ ile $\Assign{\prime}{M}$ olarak seçtiğinizde
$+$, $\times$ ve $<$ yorumları zaten belirlenir. Standart bir modelde
$\Assign{\prime}{M}$ olmaya yalnız hem !!{injective} hem !!{surjective}
olan $s\colon M\to M\setminus\{z\}$ fonksiyonları uygundur. $s$'nin
görüntü kümesi $z$'yi içeremez; aksi halde
$\lforall[x][\eq/[\Obj0][x']]$ yanlış olurdu. Bu !!{sentence},
$\Struct N$ içinde doğrudur ve dolayısıyla $\Struct M$ de onu doğru
kılmalıdır. $\Struct N$ içindeki ardıl fonksiyonu
$\Assign{\prime}{N}$ !!{injective} olduğundan ve bu olgu $\Struct N$
içinde doğru !!a{sentence} ile ifade edildiğinden $s$ de !!{injective}
olmalıdır. Ayrıca !!{surjective} olmalıdır; yoksa
$M\setminus\{z\}$ içinde $s$'nin görüntü kümesinde bulunmayan bir $x$
olur, yani $\lforall[x][(\eq[x][\Obj0]\lor
\lexists[y][\eq[y'][x]])]$ !!{sentence} ifadesi $\Struct M$ içinde
yanlış olurdu; oysa $\Struct N$ içinde doğrudur.
\end{explain}

\end{document}
